% GENERATED FILE. Edit canonical lessons, then run npm run generate:books.
% canonical-source-hash: fnv1a64:139b2e1a85aa2271
% canonical-lessons: TE-C317-iloga, TE-S174-letter-ii, TE-C317-a-taruvata, TE-C317-vacci, TE-C317-ammamma-gramam, TE-R317-a-story-in-order

\chapter{Telling a Story in Order}
\label{ch:te-telling-a-story-in-order}
\hlchaptermodality{317}

\begin{chapteropening}
\textbf{By the end of this chapter:} \emph{I can tell a short story in order with modaṭa, ā taruvāta, the -i form and civariki, and read one of about eighty words.}

Complete TE-R317-a-story-in-order: retell a day in order with modaṭa, ā taruvāta, -i and civariki, read the new letter \te{ఈ}, and say what someone did in the meantime.
\end{chapteropening}

\section[\texorpdfstring{\te{ఈలోగా} (īlōgā)}{ఈలోగా (īlōgā)}]{\te{ఈలోగా} (īlōgā) — in the meantime, before then}

\label{lesson:TE-C317-iloga}



\begin{quote}
Before the new one: say the Telugu for \textquotedblleft{}meanwhile, at that very moment\textquotedblright{}. (\emph{intalō}.) Then \textquotedblleft{}I used to read\textquotedblright{}. (\emph{cadivēvāḍini}.)
\end{quote}

\subsection*{You'll want to know: \te{ఈలోగా}}
\textbf{\te{ఈలోగా}} — \emph{īlōgā} — \textquotedblleft{}in the meantime, before then\textquotedblright{}.

\begin{quote}
\textbf{\te{అమ్మ} \te{వండుతోంది}. \te{ఈలోగా} \te{నేను} \te{ఒక} \te{పుస్తకం} \te{చదివాను}.}
\end{quote}

\begin{quote}
\emph{amma vaṇḍutōndi. īlōgā nēnu oka pustakaṁ cadivānu.} — \textquotedblleft{}Mother was cooking. In the meantime I read a book.\textquotedblright{}
\end{quote}

\begin{quote}
\textbf{\te{రైలు} \te{ఆరు} \te{గంటలకు} \te{వస్తుంది}. \te{ఈలోగా} \te{భోజనం} \te{చేద్దాం}.}
\end{quote}

\begin{quote}
\emph{railu āru gaṇṭalaku vastundi. īlōgā bhōjanaṁ cēddāṁ.} — \textquotedblleft{}The train comes at six. Before then, let's eat.\textquotedblright{}
\end{quote}

\textbf{\te{చేద్దాం}}, \textquotedblleft{}let's do\textquotedblright{}, has the ending of \textbf{\te{కలుద్దాం}}, \textquotedblleft{}let's meet\textquotedblright{}.

\textbf{\te{ఇంతలో}}, which you know, and \textbf{\te{ఈలోగా}} overlap, and you will hear both. \textbf{\te{ఇంతలో}} leans toward \textquotedblleft{}at that very moment\textquotedblright{}: something broke into the middle of a scene. \textbf{\te{ఈలోగా}} leans toward the time before something else arrives, and what someone did with it.

Its first letter is one you have not written yet. The next lesson takes it out of this word and puts it in your hand.

\begin{cousinweb}
\textbf{\te{ఈలోగా}} (\emph{īlōgā}) is native Dravidian, and it reads as \textquotedblleft{}within this\textquotedblright{}. \textbf{\te{ఈ}}, \emph{ī}, is \textquotedblleft{}this\textquotedblright{} in front of a noun, on the near-pointing base that also gives \textbf{\te{ఇక్కడ}} and \textbf{\te{ఇప్పుడు}} (DEDR 410). \textbf{\te{లోగా}}, \emph{lōgā}, \textquotedblleft{}within\textquotedblright{}, is listed in DEDR 698, in an older spelling, beside \emph{lō}, the \textbf{\te{లో}} of \textbf{\te{ఇంట్లో}} and \textbf{\te{బాల్యంలో}}.
\end{cousinweb}

\subsection*{Your turn}
Say these aloud:

\begin{itemize}
\raggedright
  \item \emph{īlōgā}, slowly, then at speed
  \item \textquotedblleft{}Mother was cooking. In the meantime I read a book.\textquotedblright{} — \emph{amma vaṇḍutōndi. īlōgā nēnu oka pustakaṁ cadivānu.}
  \item \textquotedblleft{}The train comes at six. Before then, let's eat.\textquotedblright{} — \emph{railu āru gaṇṭalaku vastundi. īlōgā bhōjanaṁ cēddāṁ.}
  \item the other word, for 'at that very moment' — \emph{intalō}
\end{itemize}

\begin{tcolorbox}[breakable,colback=teal!4,colframe=teal!35!black,title={Before you move on}]
What does \textbf{\te{ఈలోగా}} mean? (\textbf{In the meantime, before then.}) Taken apart? (\textquotedblleft{}This\textquotedblright{} + \textquotedblleft{}within\textquotedblright{}.) Which word leans toward \textquotedblleft{}at that very moment\textquotedblright{}? (\textbf{\te{ఇంతలో}}, \emph{intalō}.)
\end{tcolorbox}
\section[\texorpdfstring{\te{ఈ} (ī)}{ఈ (ī)}]{\te{ఈ} (ī) — the letter that opens \te{ఈలోగా}}

\label{lesson:TE-S174-letter-ii}



\begin{quote}
Before the new one: write \textbf{\te{ఇ}} and \textbf{\te{ీ}} once each, and name them. (Short \emph{i} at the start of a word; long \emph{ī} after a consonant, as in \textbf{\te{నీళ్ళు}}.)

Now say \textbf{\te{ఈలోగా}} (\emph{īlōgā}), \textquotedblleft{}in the meantime\textquotedblright{}. You met it a moment ago. This lesson takes its first letter out and puts it in your hand.
\end{quote}

\subsection*{Script you'll notice: \te{ఈ}}
\textbf{\te{ఈ}} — \emph{ī}, a long \emph{ī} at the start of a word.

After a consonant, long \emph{ī} is the sign \textbf{\te{ీ}}: \textbf{\te{నీళ్ళు}}, \emph{nīḷḷu}. At the start of a word there is no consonant to carry a sign, so the vowel stands as a letter of its own, \textbf{\te{ఈ}}.

\textbf{\te{ఈ}} is not \textbf{\te{ఇ}} with a mark added: it has a shape of its own. In print it is a round body with a straight bar across it, a short hook rising on the left, a long stroke sweeping up to the right, and a small curl standing beside it on the right.

On its own, \textbf{\te{ఈ}} is also a word: \textquotedblleft{}this\textquotedblright{}, in front of a noun. \textbf{\te{ఈ} \te{పుస్తకం}} is \textquotedblleft{}this book\textquotedblright{}; \textbf{\te{ఆ} \te{పుస్తకం}} is \textquotedblleft{}that book\textquotedblright{}.

\subsection*{Writing: \te{ఈ} — copy the printed shape}
Put your pen beside \textbf{\te{ఈ}} and copy it, slowly, and larger than it is printed. Keep the bar straight and the little curl on the right separate from the body.

\begin{quote}
This book does not yet say where the pen starts on this shape, or which way it travels. That is taught with real variation from school to school, and it waits here until it can be written down from a source. Copying the printed shape needs no such source, and it puts the shape in your hand in the meantime.
\end{quote}

\subsection*{Your turn}
\begin{itemize}
\raggedright
  \item \emph{Find:} \te{ఈ} at the start of \textbf{\te{ఈలోగా}}
  \item \emph{Trace it:} \te{ఈ} three times, saying \emph{ī} as you finish each one
  \item \emph{Write it:} \te{ఇ}, then \te{ఈ}, from memory, naming each
  \item \emph{Write it:} \textbf{\te{ఈలోగా}}, then \textbf{\te{ఈ} \te{పుస్తకం}}, \textquotedblleft{}this book\textquotedblright{}
  \item \emph{Read it:} \textbf{\te{ఈలోగా}} aloud once more
\end{itemize}

\begin{tcolorbox}[breakable,colback=teal!4,colframe=teal!35!black,title={Before you move on}]
Which letter is this — \te{ఈ}? (\textbf{ī}, long \emph{ī} at the start of a word.) Which word did it come from? (\textbf{\te{ఈలోగా}}, \emph{īlōgā}, \textquotedblleft{}in the meantime\textquotedblright{}.) What does \textbf{\te{ఈ}} mean on its own? (\textbf{\textquotedblleft{}This\textquotedblright{}}, in front of a noun.)
\end{tcolorbox}
\section[\texorpdfstring{\te{ఆ} \te{తరువాత} (ā taruvāta)}{ఆ తరువాత (ā taruvāta)}]{\te{ఆ} \te{తరువాత} (ā taruvāta) — first, after that, in the end}

\label{lesson:TE-C317-a-taruvata}



\begin{quote}
Say \textquotedblleft{}after, later\textquotedblright{} (\emph{taruvāta}) and \textquotedblleft{}before\textquotedblright{} (\emph{mundu}). Then the word that opened the last lesson, \textquotedblleft{}in the meantime\textquotedblright{}. (\emph{īlōgā}.)

\begin{itemize}
\raggedright
  \item \emph{Write it:} \textbf{\te{ఈ}}, once
\end{itemize}
\end{quote}

\subsection*{You'll want to know: \te{మొదట} … \te{ఆ} \te{తరువాత}}
\begin{quote}
\textbf{\te{మొదట} \te{నేను} \te{టీ} \te{తాగాను}. \te{ఆ} \te{తరువాత} \te{బడికి} \te{వెళ్ళాను}.}
\end{quote}

\begin{quote}
\emph{modaṭa nēnu ṭī tāgānu. ā taruvāta baḍiki veḷḷānu.} — \textquotedblleft{}First I drank tea. After that I went to school.\textquotedblright{}
\end{quote}

\textbf{\te{మొదట}} (\emph{modaṭa}, \textquotedblleft{}first, at first\textquotedblright{}) sits beside \textbf{\te{మొదటి}}, \textquotedblleft{}first\textquotedblright{}, which you know: \textbf{\te{మొదటి}} goes in front of a noun, \textbf{\te{మొదట}} opens a sentence. \textbf{\te{ఆ} \te{తరువాత}} is \textbf{\te{తరువాత}} with \textbf{\te{ఆ}}, \textquotedblleft{}that\textquotedblright{}, in front: \textquotedblleft{}after that\textquotedblright{}. After a time word \textbf{\te{తరువాత}} means \textquotedblleft{}after, later\textquotedblright{}:

\begin{quote}
\textbf{\te{రెండు} \te{రోజుల} \te{తరువాత} \te{అతను} \te{వచ్చాడు}.}
\end{quote}

\begin{quote}
\emph{reṇḍu rōjula taruvāta atanu vaccāḍu.} — \textquotedblleft{}Two days later he came.\textquotedblright{}
\end{quote}

\subsection*{You'll want to know: \te{చివరికి}}
\begin{quote}
\textbf{\te{చివరికి} \te{రాత్రి} \te{పది} \te{గంటలకు} \te{ఇంటికి} \te{చేరాను}.}
\end{quote}

\begin{quote}
\emph{civariki rātri padi gaṇṭalaku iṇṭiki cērānu.} — \textquotedblleft{}In the end I reached home at ten at night.\textquotedblright{}
\end{quote}

\textbf{\te{చివరికి}} is \textbf{\te{చివర}}, \textquotedblleft{}the end\textquotedblright{}, with \textbf{-\te{కి}}, \textquotedblleft{}to\textquotedblright{}: \textquotedblleft{}to the end\textquotedblright{}, so \textquotedblleft{}finally, in the end\textquotedblright{}. A told story often runs on three posts: \textbf{\te{మొదట}} … \textbf{\te{ఆ} \te{తరువాత}} … \textbf{\te{చివరికి}}.

\subsection*{Your turn}
Say these aloud:

\begin{itemize}
\raggedright
  \item \textquotedblleft{}First I drank tea. After that I went to school.\textquotedblright{} — \emph{modaṭa nēnu ṭī tāgānu. ā taruvāta baḍiki veḷḷānu.}
  \item \textquotedblleft{}Two days later he came\textquotedblright{} — \emph{reṇḍu rōjula taruvāta atanu vaccāḍu}
  \item \textquotedblleft{}In the end I reached home at ten at night\textquotedblright{} — \emph{civariki rātri padi gaṇṭalaku iṇṭiki cērānu}
  \item yesterday in three steps: first, after that, in the end
\end{itemize}

\begin{tcolorbox}[breakable,colback=teal!4,colframe=teal!35!black,title={Before you move on}]
\textquotedblleft{}First\textquotedblright{}, opening a sentence? (\textbf{\te{మొదట}}.) \textquotedblleft{}After that\textquotedblright{}? (\textbf{\te{ఆ} \te{తరువాత}}.) \textquotedblleft{}In the end\textquotedblright{}? (\textbf{\te{చివరికి}}.) What is \textbf{-\te{కి}} doing in \textbf{\te{చివరికి}}? (\textbf{\textquotedblleft{}To\textquotedblright{}: to the end.})
\end{tcolorbox}
\section[\texorpdfstring{\te{వచ్చి} (vacci)}{వచ్చి (vacci)}]{\te{వచ్చి} (vacci) — having come, and then}

\label{lesson:TE-C317-vacci}



\begin{quote}
Say \textquotedblleft{}First I drank tea. After that I went to school.\textquotedblright{} (\emph{modaṭa nēnu ṭī tāgānu. ā taruvāta baḍiki veḷḷānu.})

Two sentences, two full verbs. Telugu can fold them into one.
\end{quote}

\subsection*{You'll want to know: the -\te{ఇ} form}
\begin{quote}
\textbf{\te{నేను} \te{ఇంటికి} \te{వచ్చి} \te{భోజనం} \te{చేశాను}.}
\end{quote}

\begin{quote}
\emph{nēnu iṇṭiki vacci bhōjanaṁ cēśānu.} — \textquotedblleft{}I came home and had a meal.\textquotedblright{}
\end{quote}

The first verb ends in \textbf{-\te{ఇ}} and carries no person: \textbf{\te{వచ్చి}}, \textquotedblleft{}having come\textquotedblright{}. Only the last verb says who and when. So a chain of \textbf{-\te{ఇ}} forms runs on to one full verb at the end:

\begin{quote}
\textbf{\te{అతను} \te{అంగడికి} \te{వెళ్ళి}, \te{ఒక} \te{పుస్తకం} \te{కొని}, \te{ఇంటికి} \te{వచ్చాడు}.}
\end{quote}

\begin{quote}
\emph{atanu aṅgaḍiki veḷḷi, oka pustakaṁ koni, iṇṭiki vaccāḍu.} — \textquotedblleft{}He went to the shop, bought a book and came home.\textquotedblright{}
\end{quote}

\textbf{\te{కొని}} (\emph{koni}, \textquotedblleft{}having bought\textquotedblright{}) is the \textbf{-\te{ఇ}} form of \textbf{\te{కొను}}, \textquotedblleft{}to buy\textquotedblright{}. Many of these sit close to the past stem: \textbf{\te{వెళ్ళి}}, \textbf{\te{చదివి}}, \textbf{\te{తాగి}}; and \textbf{\te{తిను}} $\to$ \textbf{\te{తిని}}, \textbf{\te{చూడు}} $\to$ \textbf{\te{చూసి}}, \textbf{\te{చేయు}} $\to$ \textbf{\te{చేసి}}.

\subsection*{You'll want to know: -\te{ఇన} \te{తరువాత}, after doing}
\begin{quote}
\textbf{\te{భోజనం} \te{చేసిన} \te{తరువాత} \te{నేను} \te{నిద్రపోయాను}.}
\end{quote}

\begin{quote}
\emph{bhōjanaṁ cēsina taruvāta nēnu nidrapōyānu.} — \textquotedblleft{}After having the meal, I slept.\textquotedblright{}
\end{quote}

\textbf{\te{చేసిన}} is \textquotedblleft{}that was done\textquotedblright{}, and before \textbf{\te{తరువాత}} it gives \textquotedblleft{}after doing\textquotedblright{}: \textbf{\te{వచ్చిన} \te{తరువాత}}, \textquotedblleft{}after coming\textquotedblright{}; \textbf{\te{చూసిన} \te{తరువాత}}, \textquotedblleft{}after seeing\textquotedblright{}. Where \textbf{-\te{ఇ}} ties two actions tightly, \textbf{-\te{ఇన} \te{తరువాత}} marks a clear step in time.

\subsection*{Your turn}
Say these aloud:

\begin{itemize}
\raggedright
  \item \textquotedblleft{}I came home and had a meal\textquotedblright{} — \emph{nēnu iṇṭiki vacci bhōjanaṁ cēśānu}
  \item \textquotedblleft{}He went to the shop, bought a book and came home\textquotedblright{} — \emph{atanu aṅgaḍiki veḷḷi, oka pustakaṁ koni, iṇṭiki vaccāḍu}
  \item \textquotedblleft{}After having the meal, I slept\textquotedblright{} — \emph{bhōjanaṁ cēsina taruvāta nēnu nidrapōyānu}
  \item one morning of yours as one sentence, with two -\te{ఇ} forms and one full verb
\end{itemize}

\begin{tcolorbox}[breakable,colback=teal!4,colframe=teal!35!black,title={Before you move on}]
\textquotedblleft{}Having come\textquotedblright{}? (\textbf{\te{వచ్చి}}.) \textquotedblleft{}Having eaten\textquotedblright{}? (\textbf{\te{తిని}}.) Which verb in a chain says who and when? (\textbf{The last.}) \textquotedblleft{}After seeing\textquotedblright{}? (\textbf{\te{చూసిన} \te{తరువాత}}.)
\end{tcolorbox}
\section[(reading)]{(reading) — a visit to Grandmother's village}

\label{lesson:TE-C317-ammamma-gramam}



\begin{quote}
About eighty words, all met before; the few new forms are glossed after it. The writer is a man.
\end{quote}

\begin{tcolorbox}[breakable,skin=enhanced,colback=orange!6,colframe=orange!45!black,boxrule=0.5pt,arc=1mm,left=6pt,right=6pt,top=4pt,bottom=4pt,fonttitle=\bfseries,title={Read this}]
\begin{quote}
\te{ఇటీవల} \te{నేను} \te{మా} \te{అమ్మమ్మ} \te{గ్రామానికి} \te{వెళ్ళాను}. \te{మొదట} \te{నేను} \te{బస్సులో} \te{కూర్చున్నాను}. \te{బస్సు} \te{వెళ్తున్నప్పుడు} \te{నేను} \te{ఒక} \te{పుస్తకం} \te{చదువుతున్నాను}. \te{అకస్మాత్తుగా} \te{వర్షం} \te{పడింది}. \te{చివరికి} \te{సాయంత్రం} \te{గ్రామానికి} \te{చేరాను}. \te{నేను} \te{బస్సు} \te{దిగి} \te{నెమ్మదిగా} \te{నడిచాను}.
\end{quote}

\begin{quote}
\te{ఈలోగా} \te{అమ్మమ్మ} \te{భోజనం} \te{వండింది}. \te{అమ్మమ్మ} \te{నన్ను} \te{చూసి} \te{నవ్వింది}. \te{మేము} \te{కలిసి} \te{భోజనం} \te{చేశాము}. \te{భోజనం} \te{చేసిన} \te{తరువాత} \te{మేము} \te{టీ} \te{తాగాము}. \te{ఆ} \te{తరువాత} \te{అమ్మమ్మ} \te{నా} \te{బాల్యం} \te{గురించి} \te{చెప్పింది}. \te{బాల్యంలో} \te{నేను} \te{ప్రతి} \te{రోజు} \te{అమ్మమ్మతో} \te{అంగడికి} \te{వెళ్ళేవాడిని}. \te{అప్పుడు} \te{గ్రామంలో} \te{ఒక} \te{చిన్న} \te{అంగడి} \te{మాత్రమే} \te{ఉండేది}. \te{ఇప్పుడు} \te{అక్కడ} \te{ఒక} \te{పెద్ద} \te{హోటలు} \te{కూడా} \te{ఉంది}.
\end{quote}

\begin{quote}
\te{రాత్రి} \te{నేను} \te{తొమ్మిది} \te{గంటలకు} \te{నిద్రపోయాను}. \te{అది} \te{చాలా} \te{మంచి} \te{రోజు}.
\end{quote}

Read it once without stopping, then once more, sorting every verb: scene, event, or habit.
\end{tcolorbox}

\begin{culture}
New forms: \textbf{\te{కూర్చున్నాను}} (\emph{kūrcunnānu}, \textquotedblleft{}I sat\textquotedblright{}), \textbf{\te{కలిసి}} (\emph{kalisi}, \textquotedblleft{}together\textquotedblright{}), \textbf{\te{నన్ను}} (\emph{nannu}, \textquotedblleft{}me\textquotedblright{}), \textbf{\te{నవ్వింది}} (\emph{navvindi}, \textquotedblleft{}she smiled\textquotedblright{}), and \textbf{-\te{తో}}, \textquotedblleft{}with\textquotedblright{}, in \textbf{\te{అమ్మమ్మతో}}.

The scene is \textbf{\te{చదువుతున్నాను}}; the habit is \textbf{\te{వెళ్ళేవాడిని}}; the background is \textbf{\te{ఉండేది}}. Every other verb is an event.
\end{culture}

\subsection*{Your turn}
Say these aloud:

\begin{itemize}
\raggedright
  \item what the writer was doing when it rained
  \item what Grandmother did in the meantime
  \item what the writer used to do every day
\end{itemize}

\begin{tcolorbox}[breakable,colback=teal!4,colframe=teal!35!black,title={Before you move on}]
Why is \textbf{\te{ఉండేది}} in the habit form? (The small shop \textbf{used to be} there: it is background.)
\end{tcolorbox}
\section[(review)]{Review — a story in order}

\label{lesson:TE-R317-a-story-in-order}



\begin{quote}
Say \textquotedblleft{}in the meantime\textquotedblright{}. (\emph{īlōgā}.)

\begin{itemize}
\raggedright
  \item \emph{Write it:} its first letter, \textbf{\te{ఈ}}, once from memory
\end{itemize}
\end{quote}

\subsection*{Your turn}
\begin{itemize}
\raggedright
  \item \emph{Say it:} First I drank tea. After that I went to the town.
  \item \emph{Say it:} I came home and had a meal. In the meantime it rained.
  \item \emph{Say it:} After seeing the film, we went home. In the end we slept.
  \item \emph{Say it:} from the reading, what the writer was doing on the bus, and what he used to do every day
  \item \emph{Write it:} \textbf{\te{ఈలోగా}}, then \textbf{\te{ఈ} \te{పుస్తకం}}, \textquotedblleft{}this book\textquotedblright{}
\end{itemize}

\begin{tcolorbox}[breakable,colback=teal!4,colframe=teal!35!black,title={Before you move on}]
\textquotedblleft{}First\textquotedblright{}? (\textbf{\te{మొదట}}.) \textquotedblleft{}After that\textquotedblright{}? (\textbf{\te{ఆ} \te{తరువాత}}.) \textquotedblleft{}In the end\textquotedblright{}? (\textbf{\te{చివరికి}}.) \textquotedblleft{}Having come\textquotedblright{}? (\textbf{\te{వచ్చి}}.) \textquotedblleft{}In the meantime\textquotedblright{}? (\textbf{\te{ఈలోగా}}, which opens with \textbf{\te{ఈ}}.) In the reading, what did Grandmother do in the meantime? (\textbf{She cooked a meal.})
\end{tcolorbox}
